\SourcePage{569}{572}
\section{Isotopic invariance}

At present there is still no complete theory of the so-called \emph{nuclear forces}: the forces acting between nuclear particles (\emph{nucleons}) and binding them together within an atomic nucleus. Consequently, present descriptions of nuclear forces must rely on experiment to a much greater degree than would be needed if a consistent theory were available.

The two types of particle classed as nucleons differ first of all in their electrical properties: protons ($p$) carry positive charge, whereas neutrons ($n$) are electrically neutral. Both nevertheless have spin $1/2$, and their masses are nearly equal (the proton and neutron masses are respectively $1836{.}1$ and $1838{.}6$ electron masses). This resemblance is not accidental. Despite their differing electrical properties, proton and neutron are very similar particles, and the similarity is of fundamental significance.

If the comparatively weak electrical forces are set aside, the interaction forces between two protons prove to be very similar to those between two neutrons. This property is called \emph{charge symmetry} of the nuclear forces\footnote{It is manifested, for example, by the similarity between the properties---binding energy, energy spectrum, and so forth---of so-called mirror nuclei, which are related by replacing every proton with a neutron and every neutron with a proton.}.

To the accuracy with which this symmetry holds, one may in particular say that a two-proton system ($pp$) and a two-neutron system ($nn$) possess states with identical properties. It is essential here that protons and neutrons obey the same statistics, namely Fermi statistics. The $pp$ and $nn$ systems consequently permit only states with the same wave-function symmetry: $\psi(\mathbf r_1,\sigma_1;\mathbf r_2,\sigma_2)$ must be antisymmetric under simultaneous interchange of the particle coordinates and spins.

\SourcePage{570}{573}
Charge symmetry, however, is only one manifestation of a still deeper physical resemblance between proton and neutron, known as \emph{isotopic invariance}\footnote{The term \emph{isobaric invariance} is also used in the literature.}. This deeper regularity produces an analogy not only between $pp$ and $nn$, which are related by interchanging every proton and neutron, but also with the $pn$ system of unlike particles. The analogy cannot of course be complete, since a $pn$ system consists of nonidentical particles and its possible states need not in every case be restricted to antisymmetric wave functions. Among its allowed states, however, there are states whose properties agree to high accuracy with those of systems formed by two identical nucleons\footnote{This was demonstrated by \emph{Breit, Condon, and Present} (O.~Breit, E.~U.~Condon, R.~D.~Present, 1936) through analysis of experimental data on neutron--proton and proton--proton scattering.}. These states are naturally described by antisymmetric wave functions; the remaining $pn$ states have symmetric wave functions and are absent from $pp$ and $nn$ systems.

Like charge symmetry, isotopic invariance is valid only when electromagnetic interaction is neglected. Its approximate character has another source in the small difference between the neutron and proton masses: exact neutron--proton symmetry would plainly require their masses to agree exactly\footnote{It is reasonable to suppose that the neutron--proton mass difference itself is in fact electromagnetic in origin.}.

A convenient formal apparatus can be introduced to describe isotopic invariance. It arises naturally upon observing that this invariance amounts to the possibility of classifying the states of a nucleon system by the symmetry of its coordinate--spin wave functions $\psi$, without regard to which of the two types the nucleons belong. The required apparatus must therefore introduce a new quantum number for characterizing the system's states, such that its value determines the symmetry of $\psi$ uniquely. We have already met an analogous situation for a system of spin-$1/2$ particles. In \S~63 we saw that specifying the total spin $S$ of such a system uniquely fixes the symmetry of

\SourcePage{571}{574}
its coordinate wave function $\varphi$, independently of which of the two possible values, $\pm1/2$, are taken by the individual spin projections $\sigma$.

It is therefore natural, in a formal account of isotopic invariance, to regard neutron and proton as two different \emph{charge states} of a single particle, the nucleon. The states are distinguished by the projection of a new vector $\boldsymbol\tau$ whose formal properties are those of a spin-$1/2$ vector. This new quantity is conventionally called \emph{isotopic spin}, or simply \emph{isospin}\footnote{It was introduced by \emph{Heisenberg} (1932) and applied to isotopic invariance by \emph{Cassen and Condon} (B.~Cassen, E.~U.~Condon, 1936).}. It is a vector in an auxiliary ``isotopic space'' $\xi\eta\zeta$, which of course has no relation to physical space.

The projection of a nucleon's isospin on the $\zeta$ axis has only the two values $\tau_\zeta=\pm1/2$. By convention, $+1/2$ is assigned to the proton and $-1/2$ to the neutron\footnote{The reverse convention is also used in the literature.}. The isospins of several nucleons combine into the system's total isospin by the ordinary rules for adding spins. The $\zeta$ component of that total isospin is the sum of the $\tau_\zeta$ values of all its constituent particles. For a nucleus containing $Z$ protons, or having atomic number $Z$, and $N$ neutrons, with mass number $A=Z+N$, one has
\begin{equation}
 T_\zeta=\sum\tau_\zeta=\frac{Z-N}{2}=Z-\frac A2,
 \tag{116.1}
\end{equation}
so that at fixed nucleon number $T_\zeta$ determines the system's total charge. Exact conservation of $T_\zeta$ is therefore evident; it merely expresses conservation of charge.

The magnitude $T$ of the total isospin determines the symmetry of the ``charge part'' $\omega$ of the system's wave function, just as the total spin $S$ determines the symmetry of its spin wave function. It thereby determines the symmetry of the coordinate--spin, or ordinary, wave function $\psi$ as well. Indeed, the total wave function of a nucleon system, the product $\psi\omega$, must have a prescribed symmetry: as for all fermions, it must be antisymmetric under simultaneous interchange of the particle coordinates, spins, and ``charge variables'' $\tau_\zeta$. Thus, in the present scheme, the definite symmetry of $\psi$ for

\SourcePage{572}{575}
any nucleon system is expressed precisely by conservation of $T$.

Equivalently, isotopic invariance means invariance of the system's properties under arbitrary rotations in isotopic space. States that differ only in $T_\zeta$, with $T$ and all other quantum numbers fixed, have the same properties. Charge symmetry---invariance under exchanging all neutrons with protons and conversely---is a special instance. It is represented here by invariance under simultaneous reversal of every $\tau_\zeta$, that is, under a rotation in isospace through $180^\circ$ about an axis in the $\xi\eta$ plane.

The manifest violation of isotopic invariance by the Coulomb interaction is also formally apparent in this scheme: that interaction depends on charge, hence on the $\zeta$ components of isospin, which are not invariant under rotations in $\xi\eta\zeta$ space.

Consider, for example, a two-nucleon system. Its total isospin can be $T=1$ or $T=0$. For $T=1$, the allowed projections are $T_\zeta=1,0,-1$. By (116.1), these correspond to charges $2$, $1$, and $0$; the $T=1$ system may therefore be realized as $pp$, $pn$, or $nn$. Its charge wave function $\omega$ is symmetric, just as a spin value $S=1$ corresponds to a symmetric spin function (compare \S~62). Hence $T=1$ corresponds to states with antisymmetric ordinary wave functions $\psi$. For $T=0$, only $T_\zeta=0$ is possible and $\omega$ is antisymmetric; this case consequently comprises $pn$ states with symmetric wave functions $\psi$.

Isospin is represented by an operator $\hat{\boldsymbol\tau}$ acting on the charge variable $\tau_\zeta$ in a wave function, in the same manner as the spin operator $\hat{\mathbf s}$ acts on the spin variable $\sigma$. Because the formal analogy is complete, the operators $\hat\tau_\xi$, $\hat\tau_\eta$, and $\hat\tau_\zeta$ are expressed by the same Pauli matrices (55.7) as $\hat s_x$, $\hat s_y$, and $\hat s_z$.

Several combinations of these operators have an immediate interpretation. The sum
\[
 \hat\tau_+=\hat\tau_\xi+i\hat\tau_\eta=
 \begin{pmatrix}0&1\\0&0\end{pmatrix}
\]
is an operator that converts a neutron wave function into a proton wave function, while annihilating a proton

\SourcePage{573}{576}
wave function. Similarly,
\[
 \hat\tau_-=\hat\tau_\xi-i\hat\tau_\eta=
 \begin{pmatrix}0&0\\1&0\end{pmatrix}
\]
converts a proton into a neutron and annihilates a neutron. Finally, the operator
\[
 \frac12+\hat\tau_\zeta=\begin{pmatrix}1&0\\0&0\end{pmatrix}
\]
leaves a proton function unchanged and annihilates a neutron function; it may be called the nucleon charge operator, in units of $e$.

We shall also show how the interchange operator $\hat P$ for two particles can be written in terms of their isospin operators $\hat{\boldsymbol\tau}_1$ and $\hat{\boldsymbol\tau}_2$. By definition, its action on the wave function $\psi(\mathbf r_1,\sigma_1;\mathbf r_2,\sigma_2)$ of a two-particle system interchanges the particles' coordinates and spins, or the variables $\mathbf r_1,\sigma_1$ with $\mathbf r_2,\sigma_2$. Its eigenvalues are $\pm1$, obtained on symmetric and antisymmetric functions $\psi$:
\begin{equation}
 \hat P\psi_{\mathrm{sym}}=\psi_{\mathrm{sym}},\qquad
 \hat P\psi_{\mathrm{antisym}}=-\psi_{\mathrm{antisym}}.
 \tag{116.2}
\end{equation}

As shown above, the functions $\psi_{\mathrm{sym}}$ and $\psi_{\mathrm{antisym}}$ correspond respectively to charge functions $\omega_T$ with total isospin $T=0$ and $T=1$. Thus, if $\hat P$ is to be represented in a form acting on the charge variables, it must satisfy
\begin{equation}
 \hat P\omega_0=\omega_0,\qquad \hat P\omega_1=-\omega_1.
 \tag{116.3}
\end{equation}
The operator $1-\hat{\mathbf T}^2$ has these properties, as follows at once from the fact that $\omega_T$ is an eigenfunction of $\hat{\mathbf T}^2$ with eigenvalue $T(T+1)$. Finally, writing $\hat{\mathbf T}=\hat{\boldsymbol\tau}_1+\hat{\boldsymbol\tau}_2$ and using the fixed equal values $\tau(\tau+1)=3/4$ of $\hat{\boldsymbol\tau}_1^2$ and $\hat{\boldsymbol\tau}_2^2$, we obtain the required final expression\footnote{An operator of this form, constructed from the ordinary particle spins, was encountered in the problems to \S~62.}:
\begin{equation}
 \hat P=1-\hat{\mathbf T}^2=-\frac12
 -2\hat{\boldsymbol\tau}_1\hat{\boldsymbol\tau}_2.
 \tag{116.4}
\end{equation}

\SourcePage{574}{577}
Matrix elements of various physical quantities of a nucleon system obey definite selection rules in isospin (L.~A.~Radicati, 1952). Let $F$ be a quantity of any tensor character that is additive, in the sense that its value for the system is the sum of its values for the separate nucleons. Write its operator as
\[
 \hat F=\sum_p\hat f_p+\sum_n\hat f_n,
\]
where the sums extend over all protons and all neutrons in the system. This expression may be rewritten identically as
\begin{equation}
 \begin{split}
 \hat F&=\sum\left(\frac12+\hat\tau_\zeta\right)\hat f_p
 +\sum\left(\frac12-\hat\tau_\zeta\right)\hat f_n\\
 &=\frac12\sum(\hat f_p+\hat f_n)
 +\sum(\hat f_p-\hat f_n)\hat\tau_\zeta,
 \end{split}
 \tag{116.5}
\end{equation}
where every sum in each term now extends over all nucleons, both protons and neutrons. The first term in (116.5) is a scalar in isospace, while the second is the $\zeta$ component of a vector in that space. The corresponding isospin selection rules are therefore the same as the orbital-angular-momentum rules for scalars and vectors in ordinary space (see \S~29). An isotopic scalar permits only transitions with unchanged $T$. The $\zeta$ component of an isotopic vector has matrix elements only for $\Delta T=0,\pm1$; in addition, transitions with $\Delta T=0$ between states with $T_\zeta=0$ are forbidden, namely for systems containing equal numbers of neutrons and protons. The latter rule follows because the matrix element for a transition with $\Delta T=0$ is proportional to $T_\zeta$; see (29.7).

For the nuclear dipole moment, for instance, $f_p=e\mathbf r$ and $f_n=0$. The first term in (116.5) is then
\[
 \frac e2\sum\mathbf r=\frac{e}{2m}\sum m\mathbf r,
\]
which is proportional to the radius vector of the centre of mass and can be made zero by a suitable choice of origin. In other words, the nuclear dipole moment reduces to the $\zeta$ component of an isotopic vector.
